%% parts/p03/maps-elementarity.tex --- v1.1 Part IV, lines 6049--6509.
%% The six source-complete results retain proved status and point, in source
%% order, to annex:p3.maps.1--annex:p3.maps.6.
%% Source-accounting note: the source-local evidence convention is superseded
%% by the handbook's compile-enforced status grammar.  The three source
%% exercises at 6422--6430 and 6494--6508 are deliberately retired as exercise
%% apparatus: their diagram/parameter-expansion content is subsumed by the
%% exact entries below, and the third merely requests a counterexample to the
%% retained coherent-chain warning.  All three source examples and both warnings
%% are retained in compressed form.

\chapter{Maps, Substructures, and Elementarity}
\label{sec:p03-maps-elementarity}

Throughout this part, the ambient metatheory includes the Axiom of Choice, a
signature is single-sorted and first-order with equality, and every structure
has nonempty domain.  The map definitions extend to many-sorted signatures by
requiring sort preservation.  Write \(M=\dom\mathcal M\) and
\(N=\dom\mathcal N\) for the domains of two \(\Lang\)-structures.  Standard
references are \cite{ChangKeisler1990,Hodges1993,Marker2002}.  The source
chapter's three evidence descriptions are implemented here by the handbook's
compile-enforced \texttt{proved}, \texttt{sketched}, and \texttt{cited}
statuses rather than repeated as a separate mathematical convention.

\section{Homomorphisms and isomorphisms}
\label{sec:p03-map-notions}

\begin{hbdefinition}{def:p03-homomorphism}{Homomorphism}
A \emph{homomorphism} \(h:\mathcal M\to\mathcal N\) is a function \(h:M\to N\)
such that, for every constant symbol \(c\), every \(n\)-ary function symbol
\(f\), every \(n\)-ary relation symbol \(R\), and every \(\bar a\in M^n\),
\[
 h(c^{\mathcal M})=c^{\mathcal N},\qquad
 h\bigl(f^{\mathcal M}(\bar a)\bigr)=f^{\mathcal N}(h\bar a),
\]
and
\[
 \mathcal M\entails R(\bar a)
 \quad\Longrightarrow\quad
 \mathcal N\entails R(h\bar a).
\]
Here \(h\bar a\) is the coordinatewise image of the tuple.  The relation
clause is only forward.
\end{hbdefinition}

\begin{hbdefinition}{def:p03-embedding}{Embedding}
An \emph{embedding} \(h:\mathcal M\to\mathcal N\) is an injective
homomorphism that also reflects every relation symbol:
\[
 \mathcal M\entails R(\bar a)
 \quad\Longleftrightarrow\quad
 \mathcal N\entails R(h\bar a).
\]
Because equality is logical identity, injectivity is precisely reflection of
equality.
\end{hbdefinition}

\begin{hbdefinition}{def:p03-isomorphism}{Isomorphism and automorphism}
An \emph{isomorphism} is a surjective embedding.  An \emph{automorphism} of
\(\mathcal M\) is an isomorphism \(\mathcal M\to\mathcal M\).
\index{automorphism}
\end{hbdefinition}

\begin{hblemma}{thm:p03-term-preservation}{proved=annex:p3.maps.1}{Preservation of terms}
If \(h:\mathcal M\to\mathcal N\) is a homomorphism, \(t(\bar x)\) is a
\(\Lang\)-term, and \(\bar a\) is a tuple from \(M\) of the displayed arity,
then
\[
 h\bigl(t^{\mathcal M}(\bar a)\bigr)=t^{\mathcal N}(h\bar a).
\]
\end{hblemma}

\begin{hbwarning}{def:p03-homomorphism-limit}{Limits of the map notions}
A homomorphism need not preserve negated atomic formulas, and an embedding
need not preserve quantified formulas.  Preservation of all first-order
formulas is the additional elementarity condition defined below.
\end{hbwarning}

\section{Substructures and images}
\label{sec:p03-substructures}

\begin{hbdefinition}{def:p03-substructure}{Substructure}
The structure \(\mathcal M\) is a \emph{substructure} of \(\mathcal N\),
written \(\mathcal M\subseteq\mathcal N\), iff \(M\subseteq N\) and the
inclusion is an embedding.  Equivalently, every constant of \(\mathcal N\)
lies in \(M\), \(M\) is closed under every function of \(\mathcal N\), and
each relation of \(\mathcal M\) is the restriction to \(M\) of the
corresponding relation of \(\mathcal N\).
\end{hbdefinition}

\begin{hbdefinition}{def:p03-generated-substructure}{Generated substructure}
Let \(A\subseteq N\).  If \(A\ne\varnothing\), or if \(\Lang\) has a constant
symbol, the \emph{substructure of \(\mathcal N\) generated by \(A\)} is the
intersection of all substructures of \(\mathcal N\) whose domains contain
\(A\).  Its domain consists exactly of the values
\(t^{\mathcal N}(\bar a)\), where \(t\) is a term and \(\bar a\) is a finite
tuple from \(A\), including values of constant terms.  Under the nonempty-domain
convention, if \(A=\varnothing\) and the signature has no constant, a least
substructure need not exist.
\end{hbdefinition}

\begin{hbexample}{def:p03-nonelementary-substructure}{A substructure that is not elementary}
In the relational signature \(\{<\}\), the inclusion
\((\N,<)\subseteq(\mathbb Z,<)\) is a substructure embedding, but with the
parameter \(0\in\N\),
\[
 (\mathbb Z,<)\entails\exists y\,(y<0),
 \qquad
 (\N,<)\not\entails\exists y\,(y<0).
\]
Thus agreement on atomic formulas does not imply agreement on quantified
formulas.
\end{hbexample}

\begin{hblemma}{thm:p03-isomorphism-invariance}{proved=annex:p3.maps.2}{Isomorphism invariance of satisfaction}
If \(h:\mathcal M\to\mathcal N\) is an isomorphism, then for every
\(\Lang\)-formula \(\varphi(x_1,\ldots,x_k)\) and every
\(\bar a\in M^k\),
\[
 \mathcal M\entails\varphi(\bar a)
 \quad\Longleftrightarrow\quad
 \mathcal N\entails\varphi(h\bar a).
\]
In particular, isomorphic structures satisfy the same sentences.
\end{hblemma}

\begin{hblemma}{thm:p03-embedding-image}{proved=annex:p3.maps.3}{Image of an embedding}
Let \(h:\mathcal M\to\mathcal N\) be an embedding.  Then \(h(M)\) contains
every \(c^{\mathcal N}\) and is closed under every \(f^{\mathcal N}\).  The
structure \(h[\mathcal M]\) with domain \(h(M)\), the constants and functions
restricted from \(\mathcal N\), and every relation restricted from
\(\mathcal N\), is a substructure of \(\mathcal N\); the corestriction
\(h:\mathcal M\to h[\mathcal M]\) is an isomorphism.
\end{hblemma}

\section{Elementary maps and the Tarski--Vaught test}
\label{sec:p03-elementary-maps}

\begin{hbdefinition}{def:p03-elementary-map}{Elementary map}
Let \(A\subseteq M\).  A map \(h:A\to N\) is \emph{elementary} iff, for every
\(\Lang\)-formula \(\varphi(x_1,\ldots,x_n)\) and every
\(a_1,\ldots,a_n\in A\),
\[
 \mathcal M\entails\varphi(a_1,\ldots,a_n)
 \quad\Longleftrightarrow\quad
 \mathcal N\entails\varphi(h(a_1),\ldots,h(a_n)).
\]
When \(A=M\), \(h\) is an \emph{elementary embedding}.
\index{elementary embedding}
\end{hbdefinition}

\begin{hbdefinition}{def:p03-elementary-substructure}{Elementary substructure}
If \(\mathcal M\subseteq\mathcal N\) and the inclusion is elementary, then
\(\mathcal M\) is an \emph{elementary substructure} of \(\mathcal N\), written
\(\mathcal M\preccurlyeq\mathcal N\); \(\mathcal N\) is then an
\emph{elementary extension} of \(\mathcal M\).
\index{elementary extension}
\end{hbdefinition}

The notation \(\mathcal M\subseteq\mathcal N\) concerns interpretations of
signature symbols; \(\mathcal M\preccurlyeq\mathcal N\) additionally asserts
agreement on every formula with parameters from \(M\).

\begin{hbtheorem}{thm:p03-tarski-vaught}{proved=annex:p3.maps.4}{Tarski--Vaught test}
Let \(\mathcal M\subseteq\mathcal N\).  Then
\(\mathcal M\preccurlyeq\mathcal N\) iff, for every formula
\(\varphi(y,\bar x)\) and every finite tuple \(\bar a\) from \(M\),
\[
 \mathcal N\entails\exists y\,\varphi(y,\bar a)
 \quad\Longrightarrow\quad
 \mathcal N\entails\varphi(b,\bar a)
 \text{ for some }b\in M.
\]
\end{hbtheorem}

\begin{hbexample}{def:p03-tarski-vaught-acf}{Applying the Tarski--Vaught test}
If \(K\subseteq L\) are algebraically closed fields in the ring language, the
quantifier-elimination theorem for algebraically closed fields yields
\(K\preccurlyeq L\).  Closure under the ring operations alone establishes only
\(K\subseteq L\); algebraic closure by itself is not a verification of the
displayed witness condition.
\end{hbexample}

\section{Parameter expansions and diagrams}
\label{sec:p03-diagrams}

\begin{hbdefinition}{def:p03-parameter-expansion}{Constants for parameters}
If \(A\subseteq M\), let
\[
 \Lang(A)=\Lang\cup\{c_a:a\in A\},
\]
 where the \(c_a\) are pairwise distinct new constants.  The
 \(\Lang(A)\)-expansion \(\mathcal M_A\) interprets
 \(c_a^{\mathcal M_A}=a\).  Thus \(\Lang(A)\) is a signature and
 \(\mathcal M_A\) an interpreted structure, not a theory.  The same notation
 is used below for diagrams and later for types over parameters.
\end{hbdefinition}

\begin{hbdefinition}{def:p03-diagrams}{Diagrams}
For a \(\Lang\)-structure \(\mathcal M\), its \emph{atomic diagram}
\(\operatorname{Diag}(\mathcal M)\) is the set of all atomic and negated
atomic \(\Lang(M)\)-sentences true in \(\mathcal M_M\).  Its
\emph{elementary diagram} is
\[
 \operatorname{ElDiag}(\mathcal M)
 =\operatorname{Th}_{\Lang(M)}(\mathcal M_M),
 \]
 the set of all \(\Lang(M)\)-sentences true in \(\mathcal M_M\).  Both objects
 are sets of sentences.  The first records the atomic information needed to
 test embeddings; the second records all formula instances needed to test
 elementary embeddings, so the two names are not interchangeable.
\index{atomic diagram}
\index{elementary diagram}
\end{hbdefinition}

\begin{hbtheorem}{thm:p03-diagram-tests}{proved=annex:p3.maps.5}{Diagram tests for embeddings}
Let \(\mathcal M\) and \(\mathcal N\) be \(\Lang\)-structures.
\begin{enumerate}
 \item There is an embedding \(h:\mathcal M\to\mathcal N\) iff
       \(\mathcal N\) has a \(\Lang(M)\)-expansion \(\mathcal N'\) such that
       \(\mathcal N'\entails\operatorname{Diag}(\mathcal M)\).
 \item There is an elementary embedding \(h:\mathcal M\to\mathcal N\) iff
       \(\mathcal N\) has a \(\Lang(M)\)-expansion \(\mathcal N'\) such that
       \(\mathcal N'\entails\operatorname{ElDiag}(\mathcal M)\).
\end{enumerate}
\end{hbtheorem}

\begin{hbexample}{def:p03-diagram-separation}{Atomic and elementary diagrams make different tests}
For \(\mathcal M=(\N,<)\subseteq(\mathbb Z,<)=\mathcal N\), interpret \(c_n\)
by the integer \(n\).  The expansion satisfies
\(\operatorname{Diag}(\mathcal M)\), as the inclusion is an embedding, but it
does not satisfy \(\operatorname{ElDiag}(\mathcal M)\): the sentence
\(\neg\exists y\,(y<c_0)\) is true in \(\mathcal M_M\) and false in the
expanded integer order.
\end{hbexample}

\section{Elementary chains}
\label{sec:p03-elementary-chains}

\begin{hbdefinition}{def:p03-elementary-chain}{Elementary chain}
Let \(\delta\) be a limit ordinal.  A sequence
\(\langle\mathcal M_i:i<\delta\rangle\) is an \emph{elementary chain} iff
\(\mathcal M_i\preccurlyeq\mathcal M_j\) whenever \(i<j<\delta\).  Its union
\(\bigcup_{i<\delta}\mathcal M_i\) has domain
\(\bigcup_{i<\delta}\dom\mathcal M_i\), with every symbol interpreted by the
compatible union of its interpretations along the chain.
\end{hbdefinition}

\begin{hbtheorem}{thm:p03-elementary-chain}{proved=annex:p3.maps.6}{Elementary chain theorem}
If \(\langle\mathcal M_i:i<\delta\rangle\) is an elementary chain and
\(\mathcal M=\bigcup_{i<\delta}\mathcal M_i\), then
\(\mathcal M_i\preccurlyeq\mathcal M\) for every \(i<\delta\).
\end{hbtheorem}

\begin{hbwarning}{def:p03-chain-warning}{A frequent chain error}
The union of an increasing chain of substructures is a substructure, but need
not be an elementary extension of every stage.  The elementary-chain theorem
uses elementarity of every transition inclusion, not merely coherence of the
underlying interpretations.
\end{hbwarning}

\section{Model-theory notation lookup}
\label{sec:p03-model-symbols}

Table~\ref{tab:p03-model-symbols} realizes apparatus item A-162.  It records
the registry's existing assignments; in particular, \(\cong\), never the
Kleene-equality symbol reserved for Part~4, denotes isomorphism here.
\index{elementary equivalence}

\begingroup\small
\begin{longtable}{@{}p{0.29\textwidth}p{0.61\textwidth}@{}}
\caption{Model-theory symbol subtable (apparatus A-162).}
\label{tab:p03-model-symbols}\\
\toprule
Notation & Registered meaning \\
\midrule
\endfirsthead
\toprule
Notation & Registered meaning \\
\midrule
\endhead
\(\mathcal M\preccurlyeq\mathcal N\) &
  \(\mathcal M\) is an elementary substructure of \(\mathcal N\). \\
\addlinespace
\(\mathcal M\equiv\mathcal N\) &
  \(\mathcal M\) and \(\mathcal N\) satisfy exactly the same
  \(\Lang\)-sentences (elementary equivalence). \\
\addlinespace
\(\operatorname{Diag}(\mathcal M)\) &
  Atomic and negated atomic \(\Lang(M)\)-sentences true in \(\mathcal M_M\). \\
\addlinespace
\(\operatorname{ElDiag}(\mathcal M)\) &
  All \(\Lang(M)\)-sentences true in \(\mathcal M_M\). \\
\addlinespace
\(S_n(T)\) &
  Space of complete \(n\)-types over the empty set consistent with \(T\), with
  its Stone topology; see Definition~\ref{def:p03-type-space}. \\
\addlinespace
\(\displaystyle\prod_{i\in I}\mathcal M_i/\mathcal U\) &
  Ultraproduct of the family \(\langle\mathcal M_i:i\in I\rangle\) modulo
  the ultrafilter \(\mathcal U\); see
  Definitions~\ref{def:p03-filter} and~\ref{def:p03-ultraproduct}. \\
\addlinespace
\(\mathcal M\cong\mathcal N\) &
  \(\mathcal M\) and \(\mathcal N\) are isomorphic. \\
\bottomrule
\end{longtable}
\endgroup
